\section{Einleitung}\label{sec:einleitung}
In der heutigen digital vernetzten Welt sind Netzwerke und ihre Protokolle das Rückgrat zahlreicher kritischer Anwendungen,
von einfachen Webdiensten bis hin zu komplexen \gls{iot}-Systemen.
Eine entscheidende Komponente der Netzwerksicherheit ist das Testen und die Absicherung dieser Protokolle gegen potenzielle Schwachstellen.
Hierbei kommt die Technik des Fuzzing zum Einsatz, die durch die gezielte Überlastung eines Systems mit fehlerhaften oder
unerwarteten Eingaben mögliche Sicherheitslücken aufdecken soll.
Besonders relevant ist die Analyse von Netzwerkfuzzern, die speziell für die Sicherheitstests von Netzwerkprotokollen entwickelt wurden.
Das \gls{mqtt}-Protokoll, ein leichtgewichtiges Messaging-Protokoll für kleine Sensoren und mobile Geräte, stellt in vielen
\gls{iot}- und \gls{m2m}-Anwendungen einen wichtigen Standard dar.
Aufgrund seiner weitreichenden Nutzung und der oft sensiblen Natur der übertragenen Daten, ist es essenziell, die Sicherheit
dieses Protokolls gründlich zu prüfen.
Netzwerkfuzzer wie Pulsar, \gls{afl}Net und boofuzz bieten unterschiedliche Ansätze zur Durchführung solcher Tests.
Pulsar fokussiert sich auf eine dynamische und effektive Fehlersuche durch einen Mix aus Netzwerk- und Protokoll-basiertem
Fuzzing.
\gls{afl}Net basiert auf dem bewährten Ansatz von \gls{afl} und passt diesen für Netzwerkprotokolle an, um eine effektive
Testabdeckung zu gewährleisten.
Boofuzz -- eine Weiterentwicklung des etablierten Sulley-Fuzzers -- bietet spezifische Funktionen für das Testen von
Netzwerkdiensten und deren Protokollen.
In dieser Arbeit wird die Performance dieser Netzwerkfuzzer im Kontext des \gls{mqtt}-Protokolls detailliert analysiert.
Ziel ist es, durch einen Vergleich der Fuzzing-Techniken und deren Effektivität Erkenntnisse zu gewinnen, die über die bloße
Anwendung hinausgehen und zu einer fundierten Bewertung der getesteten Tools führen.
Die Untersuchung fokussiert sich auf folgende Forschungsfragen:
\begin{questions}
    \item Welchen Performanceunterschied gibt es unter den Netzwerkfuzzern?
    \item Inwiefern sind gefundene Bugs reproduzierbar?
\end{questions}
Durch die umfassende Betrachtung dieser Fragen soll ein fundierter Vergleich der Netzwerkfuzzer ermöglicht werden, der
sowohl ihre Stärken als auch ihre Schwächen aufzeigt und wertvolle Erkenntnisse für die Weiterentwicklung und Anwendung
solcher Tools liefert.
Die Ergebnisse dieser Arbeit werden nicht nur zur Verbesserung der Sicherheitsüberprüfung des \gls{mqtt}-Protokolls beitragen,
sondern auch allgemeine Empfehlungen für den Einsatz von Fuzzern im Netzwerkbereich bieten.

